\documentclass[aps,prd,twocolumn,10pt,nofootinbib,superscriptaddress]{revtex4-2}
\usepackage{amsmath,amssymb,graphicx,booktabs}
\usepackage[colorlinks=true,linkcolor=blue,citecolor=blue,urlcolor=blue]{hyperref}
\newcommand{\QQ}{\mathcal{Q}}
\newcommand{\YY}{\mathcal{Y}}
\begin{document}

\title{A host for the dark-energy--anchored acceleration scale:\\
bimetric MOND with a DBI khronon}

\author{Carl P.\ Zimmerman}
\affiliation{Briar Creek Tech}
\date{19 August 2026}

\begin{abstract}
We previously established three necessary conditions on any relativistic home for MOND
phenomenology carrying the acceleration scale $a_0=\kappa c\sqrt{G\rho_\Lambda}$, each proved by
a construction that satisfies its predecessors and fails the next
[10.5281/zenodo.22004372]. Here we add a locality theorem that explains the pattern of those
failures and then exhibit the first construction with no computed kill against it.
\textbf{(i) Screening must be local.} The Sun sits at $0.67$--$0.74$ of its own galaxy's MOND
radius, so the solar system and the MOND regime are the \emph{same environment}: between 1~AU
and $r_M$ the local gravitational field differs by $6.3\times10^{7}$ while the potential differs
by $1.5\times$, the dark-sector density by $2.2\times$, and the velocity dispersion not at all.
Any mechanism suppressing the MOND force in the solar system must be keyed to a quantity
differing by the required $1.2$--$3.4\times10^{4}$; only the local field qualifies. This kills
environmental phase-boundary screening outright, and closed our own single-argument ansatz
(squeezed by $4.3$--$5.7\times10^{3}$). \textbf{(ii) The scale identity is host-independent.}
The promotion $a_0^2=\kappa^2G(-K(\QQ))$ --- the MOND scale \emph{is} the dark sector's
pressure --- and its derived evolution $a_0(a)/a_0(0)=(1+\sigma^2)^{-1/4}$,
$\sigma\propto\nu_0/a^3$, require only a shift-symmetric clock field on FRW. We re-derive both
with no reference to any aether: $a_0(\mathrm{rec})/a_0(0)=0.0060$ and the bound
$\nu_0\le2.36\times10^{-6}$ were never properties of the previous host.
\textbf{(iii) The construction.} Bimetric MOND [Milgrom, PRD \textbf{80}, 123536 (2009)] with
the $\beta=1$ DBI khronon $K(\QQ)=-\rho_\Lambda c^2\sqrt{1-(\QQ-\QQ_0)^2/\Lambda_D^2}$ as the
dark component and the promotion setting its acceleration scale. BIMOND's favoured-class
nonrelativistic limit is QUMOND, whose free function eats the \emph{local total field} ---
requirement R1 and the locality theorem satisfied by construction. Its legality condition
$dg_{\rm obs}/dg_{\rm bar}>0$ is met by the exponential interpolation everywhere (minimum
$0.949$ over 18 decades) even though that kernel's anomaly is non-monotone --- the exact shape
that is fatal in AeST is legal here, so the saturation trap's hypothesis is not met and the
1~AU anomaly is $10^{-3458.7}\,\mathrm{m\,s^{-2}}$. There is no unit-timelike vector, so the
vector-ghost mechanism that killed the $Z$-form has no counterpart (R2), and no
$\tilde G/G_N$ split is forced (R3). \textbf{What is owed}, stated plainly: a nonlinear
Boulware--Deser ghost analysis --- BIMOND's connection-difference interaction is not of
Hassan--Rosen form, so the known exemption does not automatically apply --- and a Boltzmann
computation, since the previous host's $0.01\sigma$ CMB pass does not transfer even though the
dust component is identical. ``No computed kill'' is not ``proven viable,'' and we say which.
\end{abstract}

\maketitle

\section{The filter, and the locality theorem}

Ref.~[10.5281/zenodo.22004372] left three necessary conditions: the free function must eat the
total-field gradient (R1); its Newtonian limit must not drive a kinetic coefficient negative
(R2); and neither repair may be bought with a cosmological/local Newton-constant split (R3).
Three subsequent constructions failed instructively. A single-argument ansatz
$X=(\QQ-\QQ_0)-\YY/2m$, whose quasi-static argument $-\Psi\QQ_0-\YY/2m$ carries the Newtonian
\emph{potential}, breaks the saturation theorem's hypothesis --- but MOND phenomenology needs
the $\YY$-term of $X$ to dominate while breaking saturation needs the $\Psi$-term to dominate,
and between 1~AU and a galaxy's MOND radius the potential supplies a contrast of only
$8.0\times$ ($7.3\times$ alt) against a required $3.4$--$4.1\times10^{4}$: squeezed out by
$4.3$--$5.7\times10^{3}$, with $m$ cancelling from the ratio. An environmental phase boundary
--- the MOND force existing only in a condensed phase --- fails more fundamentally:

\emph{Theorem (screening must be local).} The Sun sits at $8.2$~kpc while the MOND radius of a
$10^{11}M_\odot$ galaxy is $12.2$~kpc ($11.1$ alt): the solar system lies at $0.67$--$0.74\,r_M$
\emph{inside} the MOND regime's own environment. Between 1~AU and $r_M$: the local field differs
by $6.3\times10^{7}$; the total potential by $1.51$; the dark-sector density by $2.2$; the
velocity dispersion by $1.0$. Any solar-system screening must be keyed to a quantity differing
by at least the ephemeris/RAR contrast $1.2$--$3.4\times10^{4}$. Only the local field qualifies,
and it does so because of the Sun itself, not the environment. \emph{Any viable screening is a
function of the local gravitational field} --- which is why R1 demands the gradient: the
requirement is forced by the Sun's position in its own galaxy, not by formalism.

This explains the whole pattern: the $Z$-form used the local field and died on R2; the
$X$-ansatz used the potential and died on the contrast; the phase boundary used the environment
and cannot separate the regions at all. Three deaths, one cause. The constructive residue,
with the earlier result that invariants reaching the free function must contain the aether
\emph{algebraically} rather than differentially (else they enter the vector kinetic term):
a viable host eats the local field without a vector to poison.

\section{The scale identity is host-independent}

The framework's central equation is the promotion
\begin{equation}
a_0^2(\QQ)=\kappa^2 G\,(-K(\QQ)),\qquad
K(\QQ)=-M^4\sqrt{1-\tfrac{(\QQ-\QQ_0)^2}{\Lambda_D^2}},
\label{eq:promotion}
\end{equation}
$M^4=\rho_\Lambda c^2$: the MOND scale \emph{is} the dark sector's pressure, giving
$a_0=\kappa c\sqrt{G\rho_\Lambda}=9.3619\times10^{-11}\,\mathrm{m\,s^{-2}}$ at the minimum
($1.1279\times10^{-10}$ on the alternative footing; $\kappa=\tfrac12$ is \emph{fitted},
$0.529\pm0.034$). On any FRW background a shift-symmetric clock field obeys
$dK/d\QQ=I_0/a^3$, whence $(\QQ-\QQ_0)/\Lambda_D=\sigma/\sqrt{1+\sigma^2}$ with
$\sigma\propto\nu_0/a^3$ and
\begin{equation}
a_0(a)/a_0(0)=(1+\sigma^2)^{-1/4},
\label{eq:a0z}
\end{equation}
verified here by computer algebra using \emph{only} shift symmetry and FRW --- no aether, no
vector sector, no host-specific structure. Everything the framework has banked from
Eq.~\eqref{eq:a0z} --- $a_0(\mathrm{rec})/a_0(0)=0.0060$, hence MOND \emph{off} at recombination;
the environmental bound $\nu_0\le2.36\times10^{-6}$; the locality of $a_0$ --- carries over
verbatim to any host supplying a clock field. They were $K(\QQ)$ results, not AeST results.

\section{The construction}

Take bimetric MOND \cite{BIMOND} for the gravity sector: two metrics, matter coupled to
$g_{\mu\nu}$ alone, interaction built from $C^\alpha{}_{\beta\gamma}=
\Gamma^\alpha{}_{\beta\gamma}-\hat\Gamma^\alpha{}_{\beta\gamma}$ --- a tensor, with no
unit-timelike vector anywhere. Add the khronon with $\QQ=\sqrt{-g^{\mu\nu}
\partial_\mu\phi\partial_\nu\phi}$ (algebraic in $\partial\phi$) and kernel
\eqref{eq:promotion} as the dark component: $w=-1$ exactly at the minimum, dust excitations,
conserved shift charge carrying the dark abundance --- unchanged from the previous host. Set
BIMOND's acceleration scale by Eq.~\eqref{eq:promotion}.

\emph{R1 and the locality theorem: satisfied by construction.} The favoured-class
nonrelativistic limit of BIMOND is QUMOND \cite{QUMOND}: $\nabla^2\Psi=\nabla\!\cdot\!
[\nu(|\nabla\Psi_N|/a_0)\nabla\Psi_N]$, the free function eating the \emph{local total field}.
Its stability/legality condition is $dg_{\rm obs}/dg_{\rm bar}>0$. For the exponential kernel
$\nu=1/(1-e^{-\sqrt y})$ we find $\min\,dg_{\rm obs}/dg_{\rm bar}=0.949$ over eighteen decades
--- legal --- even though its anomaly $U(y)=(\nu-1)y$ peaks at $0.6476$ at $y=2.540$ and falls:
the exact non-monotonicity that is fatal in AeST. The saturation theorem's hypothesis (a fixed
curve $U(y)$) is simply not met, and the 1~AU anomaly is
$10^{-3458.7}$ ($10^{-3151.3}$ alt) $\mathrm{m\,s^{-2}}$: the $1.2$--$3.4\times10^{4}$
ephemeris/RAR gap of the AeST $\YY$-form does not exist in this host.

\emph{R2: the mechanism is absent.} No vector, no $F_{\mu\nu}F^{\mu\nu}$ term, no
$C_V=K_B-(2-K_B)J_Z$, nothing for the Newtonian limit to poison; the khronon's $\QQ$ carries no
derivative of any vector. \emph{R3:} matter sees one Newton's constant; nothing forces a split.
$\gamma_{\rm PPN}=1$ is the literature's standing statement for the favoured class (quoted, not
re-derived here).

\section{What is owed, and what is claimed}

\emph{Owed, theory side:} a nonlinear ghost analysis. ``No vector'' is not ``no ghost'': generic
bimetric interactions carry the Boulware--Deser mode \cite{BD}, the known exemption
\cite{HR} is built from $\sqrt{g^{-1}\hat g}$, and BIMOND's connection-difference interaction is
not of that form. Milgrom has argued the MOND limit is safe; we have not verified a full
analysis exists, and BIMOND must not be quoted as ghost-free. This is the R2-analog for this
host and it is open. \emph{Owed, data side:} a Boltzmann computation. The previous host's
$0.01\sigma$ CMB pass does not transfer --- the dust component is identical but the gravity
sector is not. \emph{Also unchanged:} whether the dark sector's dust remains bound in galaxies
(open problem 2d), and clusters.

\emph{Claimed:} the locality theorem; the host-independence of
Eqs.~\eqref{eq:promotion}--\eqref{eq:a0z}; and that BIMOND + DBI khronon + promotion passes
every axis the filter can currently compute --- the first construction in this programme with no
computed kill against it. \emph{Not claimed:} viability. Two named computations decide it.
$\kappa=\tfrac12$ remains fitted, never derived. There is no dark-matter \emph{particle} here;
the dark abundance is a field's conserved charge, and both the $w=-1$ result and the CMB depend
on its being present in full.

\begin{acknowledgments}
Every number is reproduced by committed self-checking scripts at
\url{https://github.com/carlzimmerman/zimmerman-formula} (\texttt{superfluid\_2026/sf01--sf07}),
each printing numbered checks and exiting non-zero on failure. Errors made and withdrawn during
this programme --- including two by its own tooling in opposite directions within one day ---
are recorded, dated, with directions stated, in \texttt{RETRACTIONS.md}. Prior art on the
$a_0$--$\Lambda$ coincidence: Milgrom (1983, 1999), Blanchet \& Le Tiec (2009), Pikhitsa
(2010), Klinkhamer \& Kopp (2011); the coincidence is old, the identity \eqref{eq:promotion}
and its derived evolution are this programme's. The underlying theories are Milgrom's (BIMOND,
QUMOND) and Skordis \& Z\l o\'snik's (AeST, the previous host).
\end{acknowledgments}

\begin{thebibliography}{9}
\bibitem{BIMOND} M.~Milgrom, Phys.\ Rev.\ D \textbf{80}, 123536 (2009).
\bibitem{QUMOND} M.~Milgrom, Mon.\ Not.\ R.\ Astron.\ Soc.\ \textbf{403}, 886 (2010).
\bibitem{BD} D.~G.~Boulware and S.~Deser, Phys.\ Rev.\ D \textbf{6}, 3368 (1972).
\bibitem{HR} S.~F.~Hassan and R.~A.~Rosen, JHEP \textbf{02}, 126 (2012).
\bibitem{SZ} C.~Skordis and T.~Z\l o\'snik, Phys.\ Rev.\ Lett.\ \textbf{127}, 161302 (2021).
\bibitem{MS08} M.~Milgrom and R.~H.~Sanders, Astrophys.\ J.\ \textbf{678}, 131 (2008).
\bibitem{TR} C.~P.~Zimmerman, DOI 10.5281/zenodo.22004372 (2026).
\end{thebibliography}
\end{document}
